\documentclass[11pt]{article}
\usepackage[margin=1in]{geometry}
\usepackage{amsmath,amssymb,bm}
\usepackage[T1]{fontenc}
\usepackage[utf8]{inputenc}
\usepackage{lmodern}
\usepackage{booktabs}
\usepackage{longtable}
\usepackage{array}
\usepackage{enumitem}
\usepackage{hyperref}

\title{Loose INS/GNSS Reference Filter Mathematical Formulation}
\author{sensor-fusion documentation}
\date{\today}

\newcommand{\R}{\mathbb{R}}
\newcommand{\T}{^{T}}
\newcommand{\crossmat}[1]{\left[#1\right]_\times}
\newcommand{\dq}{\delta q}
\newcommand{\dt}{\Delta t}

\begin{document}
\maketitle
\tableofcontents
\newpage

\section{Scope}

This document describes the loose INS/GNSS reference filter implemented in
\texttt{ekf/src/loose.rs} and \texttt{ekf/src/generated\_loose.rs}. It is used
for seeded replay, diagnostics, and comparison with the runtime augmented ESKF.
It carries ECEF navigation states, IMU bias and scale correction states, and a
residual mount quaternion for vehicle-frame non-holonomic constraints.

\section{Frames and Conventions}

\begin{longtable}{>{\raggedright\arraybackslash}p{0.12\linewidth} >{\raggedright\arraybackslash}p{0.78\linewidth}}
\toprule
\textbf{Symbol} & \textbf{Meaning} \\
\midrule
\(e\) & ECEF frame. Loose nominal position and velocity are propagated in this frame. \\
\(n\) & Local NED frame. GNSS standard deviations may be supplied in this frame and are rotated into ECEF. \\
\(s\) & Seeded sensor/vehicle frame of the IMU deltas consumed by the loose filter. \\
\(c\) & Corrected vehicle frame after residual mount correction. \\
\bottomrule
\end{longtable}

Rotations follow the project convention
\begin{equation}
  x_a=C_{ab}x_b,\qquad C(q_1\otimes q_2)=C(q_1)C(q_2).
\end{equation}
The loose attitude quaternion \(q_{es}\) maps seeded-frame vectors into ECEF:
\begin{equation}
  x_e=C_{es}(q_{es})x_s .
\end{equation}
The residual mount quaternion \(q_{cs}\) maps seeded-frame velocity into the
corrected vehicle frame:
\begin{equation}
  x_c=C_{cs}(q_{cs})x_s .
\end{equation}

The local NED-to-ECEF relation is represented by
\begin{equation}
  x_n=C_{ne}x_e,\qquad x_e=C_{ne}\T x_n .
\end{equation}

\section{Nominal State}

The nominal state contains 26 scalar fields:
\begin{equation}
  x=
  \begin{bmatrix}
  q_{es} & v_e\T & p_e\T & b_g\T & b_a\T &
  s_g\T & s_a\T & q_{cs}
  \end{bmatrix}\T .
\end{equation}
The public Rust fields are:
\begin{center}
\begin{tabular}{ll}
\toprule
Fields & Meaning \\
\midrule
\texttt{q0..q3} & \(q_{es}\), seeded frame to ECEF attitude \\
\texttt{vn,ve,vd} & ECEF velocity \(v_e\), meters per second \\
\texttt{pn,pe,pd} & ECEF position \(p_e\), meters \\
\texttt{bgx..bgz} & gyro additive correction \(b_g\), radians per second \\
\texttt{bax..baz} & accelerometer additive correction \(b_a\), meters per second squared \\
\texttt{sgx..sgz} & gyro scale correction \(s_g\) \\
\texttt{sax..saz} & accelerometer scale correction \(s_a\) \\
\texttt{qcs0..qcs3} & residual seeded-frame-to-vehicle quaternion \(q_{cs}\) \\
\bottomrule
\end{tabular}
\end{center}

The code applies scale and additive correction directly to IMU rates and
specific force:
\begin{align}
  \omega_s &= s_g\odot\omega_{\mathrm{raw},s}+b_g,\\
  f_s &= s_a\odot f_{\mathrm{raw},s}+b_a.
\end{align}
Consequently, these are correction states. A physical sensor bias appears with
the opposite sign of the additive correction required to remove it.

\section{Error State}

The 24-dimensional error state order in the Rust update and generated matrices
is
\begin{equation}
  \delta x =
  \begin{bmatrix}
  \delta p_e &
  \delta v_e &
  \delta\theta_s &
  \delta b_a &
  \delta b_g &
  \delta s_a &
  \delta s_g &
  \delta\psi_{cs}
  \end{bmatrix}\T .
\end{equation}
Index ranges are:
\begin{center}
\begin{tabular}{ll}
\toprule
Indices & Error component \\
\midrule
0..2 & ECEF position error \(\delta p_e\) \\
3..5 & ECEF velocity error \(\delta v_e\) \\
6..8 & attitude small angle \(\delta\theta_s\) \\
9..11 & accelerometer-bias correction error \(\delta b_a\) \\
12..14 & gyro-bias correction error \(\delta b_g\) \\
15..17 & accelerometer-scale error \(\delta s_a\) \\
18..20 & gyro-scale error \(\delta s_g\) \\
21..23 & residual mount error \(\delta\psi_{cs}\) \\
\bottomrule
\end{tabular}
\end{center}

\section{Nominal Propagation}

The loose filter uses a two-sample IMU increment. For each half-step \(i\in
\{1,2\}\),
\begin{align}
  \omega_{s,i} &= s_g\odot\left(\frac{\Delta\theta_{s,i}}{\dt}\right)+b_g,\\
  f_{s,i} &= s_a\odot\left(\frac{\Delta v_{s,i}}{\dt}\right)+b_a .
\end{align}
The attitude derivative includes Earth rotation:
\begin{equation}
  \dot q_{es} =
  \frac{1}{2}q_{es}\otimes
  \begin{bmatrix}0&\omega_s\T\end{bmatrix}\T
  +\frac{1}{2}
  \begin{bmatrix}
  \omega_{ie} q_3 &
  \omega_{ie} q_2 &
  -\omega_{ie} q_1 &
  -\omega_{ie} q_0
  \end{bmatrix}\T .
\end{equation}
Specific force is rotated into ECEF:
\begin{equation}
  f_e=C_{es}f_s .
\end{equation}
Velocity dynamics use the J2 ECEF gravity model and Coriolis terms:
\begin{equation}
  \dot v_e =
  f_e + g_e(p_e) - 2\,\Omega_{ie}^{e}\times v_e .
\end{equation}
The implementation writes this as component terms using
\(\omega_{ie}=7.292115\times 10^{-5}\) rad/s. Position obeys
\begin{equation}
  \dot p_e=v_e .
\end{equation}
The nominal propagation uses a predictor-corrector/trapezoidal structure:
evaluate derivatives at the current state, form a temporary state after
\(\dt\), evaluate derivatives again, then average the two derivatives for the
final position, velocity, and quaternion update. The nominal bias, scale, and
mount states are constant during propagation.

\section{Error Propagation}

The generated loose model provides
\begin{equation}
  \delta x_{k+1}=F_k\delta x_k+G_kw_k,
  \qquad
  P_{k+1}^{-}=F_kP_k^{+}F_k\T+G_kQ_kG_k\T .
\end{equation}
The 21-noise vector groups are accelerometer white noise, gyro white noise,
accelerometer-bias random walk, gyro-bias random walk, accelerometer-scale
random walk, gyro-scale random walk, and mount random walk:
\begin{equation}
  w=
  \begin{bmatrix}
  n_a & n_g & n_{ba} & n_{bg} & n_{sa} & n_{sg} & n_m
  \end{bmatrix}\T .
\end{equation}
The runtime diagonal entries are
\begin{equation}
  Q_k=\operatorname{diag}\left(
  \sigma_a^2\dt I_3,\,
  \sigma_g^2\dt I_3,\,
  \sigma_{ba}^2\dt I_3,\,
  \sigma_{bg}^2\dt I_3,\,
  \sigma_{sa}^2\dt I_3,\,
  \sigma_{sg}^2\dt I_3,\,
  \sigma_m^2\dt I_3
  \right).
\end{equation}

\section{GNSS Reference Measurements}

\subsection{Position}

GNSS position is supplied in ECEF:
\begin{equation}
  z_p=p_e+\eta_p .
\end{equation}
When a horizontal accuracy \(h_{\mathrm{acc}}\) is supplied, the implementation
forms a local NED diagonal covariance
\begin{equation}
  R_n=\operatorname{diag}\left(h_{\mathrm{acc}}^2,\,
  h_{\mathrm{acc}}^2,\,(2.5h_{\mathrm{acc}})^2\right),
\end{equation}
rotates it to ECEF, then applies a Cholesky whitening transform \(T_p\). The
actual rows fused are
\begin{equation}
  r_p=T_p(z_p-\hat p_e),\qquad H_p=T_p
  \begin{bmatrix}I_3&0&\cdots&0\end{bmatrix}.
\end{equation}

\subsection{Velocity}

GNSS velocity is supplied in ECEF, with either a scalar speed accuracy or
per-axis NED velocity standard deviations. In the per-axis case, the NED
velocity covariance is rotated to ECEF and Cholesky-whitened:
\begin{equation}
  r_v=T_v(z_v-\hat v_e),\qquad H_v=T_v
  \begin{bmatrix}0&I_3&0&\cdots&0\end{bmatrix}.
\end{equation}
If only scalar speed accuracy is available, the three ECEF velocity rows use
that scalar variance directly.

\section{Non-Holonomic Constraint}

The loose NHC predicts corrected vehicle-frame velocity from ECEF velocity:
\begin{equation}
  v_s=C_{es}\T v_e,\qquad v_c=C_{cs}v_s .
\end{equation}
The lateral and vertical constraints are
\begin{equation}
  h_y=e_y\T v_c\approx 0,\qquad
  h_z=e_z\T v_c\approx 0 .
\end{equation}
The generated rows include derivatives with respect to velocity, attitude, and
residual mount. For the mount block, with left perturbation
\(q_{cs}^{true}=\dq(\delta\psi_{cs})\otimes q_{cs}\),
\begin{equation}
  \frac{\partial(e_i\T v_c)}{\partial\delta\psi_{cs}}
  =-e_i\T\crossmat{v_c},\qquad i\in\{y,z\}.
\end{equation}
NHC rows are gated by speed and a quasi-static IMU check:
\begin{equation}
  \|v_e\| \ge 0.5\ \mathrm{m/s},\qquad
  \|\omega_s\|<0.03\ \mathrm{rad/s},\qquad
  \left|\|f_s\|-9.81\right|<0.2\ \mathrm{m/s^2}.
\end{equation}
A scalar chi-square gate prevents large residuals from being fused.

\section{Batch Update and Injection}

The loose filter batches up to eight rows: three position rows, three velocity
rows, and two NHC rows. It uses a sequential covariance update in f64 shadow
storage:
\begin{align}
  S_i &= H_iPH_i\T+R_i,\\
  K_i &= PH_i\T S_i^{-1},\\
  \delta x &\leftarrow \delta x + K_i(r_i-H_i\delta x),\\
  P &\leftarrow P-K_iS_iK_i\T .
\end{align}
After the batch, the accumulated error is injected:
\begin{align}
  p_e^{+}&=p_e^{-}+\delta p_e,\qquad
  v_e^{+}=v_e^{-}+\delta v_e,\\
  q_{es}^{+}&=\dq(\delta\theta_s)\otimes q_{es}^{-},\\
  b_a^{+}&=b_a^{-}+\delta b_a,\qquad
  b_g^{+}=b_g^{-}+\delta b_g,\\
  s_a^{+}&=s_a^{-}+\delta s_a,\qquad
  s_g^{+}=s_g^{-}+\delta s_g,\\
  q_{cs}^{+}&=\dq(\delta\psi_{cs})\otimes q_{cs}^{-}.
\end{align}
Both quaternions are normalized. Unlike the runtime ESKF implementation, the
loose batch update currently does not apply an explicit post-injection reset
Jacobian to covariance.

\section{Seeded ECEF/NED Initialization}

For visualizer seeded replay, a yaw in the local NED frame is split into an
ECEF seed attitude and an identity residual mount. The utility constructs a
local NED-to-ECEF attitude from latitude and longitude, applies the yaw, and
initializes \(q_{cs}=1\). This lets the loose path consume the same
pre-rotated IMU stream used by ESKF comparisons while propagating navigation in
ECEF.

\section{Observability and Known Limitations}

\begin{itemize}[leftmargin=1.5em]
\item GNSS position and velocity directly observe only ECEF position and
velocity. Attitude, biases, scales, and mount are updated through covariance
coupling and inertial dynamics.
\item Lateral and vertical NHC rows are the direct measurement path for
residual mount. They are weak during low speed, poor seed attitude, or
high-dynamic maneuvers outside the gate.
\item Accelerometer vertical bias and accelerometer vertical scale are strongly
coupled in ground-vehicle motion dominated by gravity. The visualizer commonly
keeps accelerometer scale fixed for diagnostics.
\item The additive bias states are correction states. Diagnostic plots should
not be interpreted as physical sensor bias without applying the sign convention.
\item Mount, attitude, gyro bias, and gyro scale can become correlated during
turns. Compare seed, residual, and composed mount separately when diagnosing
mount convergence.
\end{itemize}

\end{document}
